\section{Palettes and simultaneous color returns}\label{sec:palettes}

Throughout this section and Section~\ref{sec:recurrence}, $D$ is a division
algebra of odd reduced degree $n\geq3$.  Suppose that the second alternative
of Proposition~\ref{quo:dichotomy} occurs: there is a homomorphism
\[
 \phi:U(k)\longrightarrow\mathcal F,
 \qquad \phi(V_0)=\mathcal F,\qquad R\subseteq\ker\phi,
\]
where $\mathcal F$ is finite, nonabelian, and simple.  We call the values of
$\phi$ \emph{colors}.  Every central unitary scalar has color $1$, since
its image centralizes the surjective image $\mathcal F$.

We shall rule out this homomorphism by comparing multiplicative colors
with additive translations on the Hermitian elements of $D$.  Two facts
make the comparison possible.  First, a probability law on color
configurations will retain every color while being invariant under all
additive translations.  Second, the following finite-group obstruction
will prohibit certain simultaneous returns of those colors.

The following finite-group theorem provides the distinction between colors
that the recurrence argument will preserve. Its double-coset assertion is
\cite[Lemma~5.2 and Section~6]{MPNF}. We include the proof, together with
the additional assertion about abelian subgroups, in
Appendix~\ref{sec:finite-groups}.

\begin{theorem}[Finite simple-group obstruction]\label{fin:obstruction}
Let $F$ be a finite nonabelian simple group. There is a nonempty proper
conjugacy-invariant subset $\mathcal S\subset F$, with $1\notin\mathcal S$,
having the following properties.
\begin{enumerate}[label=\textup{(\roman*)}]
\item If $W\in\mathcal S$ and $Z\notin\mathcal S$ commute, then there
exists $B\in F$ such that
\[
 ZW\notin C_F(B)(ZW^{-1})C_F(BZ^2).
\]
\item For every abelian subgroup $A\subset F$ with
$A\cap\mathcal S\ne\varnothing$, the complement $A\setminus\mathcal S$
is a subgroup of $A$.
\end{enumerate}
In particular, $\mathcal S$ is invariant under inversion.
\end{theorem}


Fix a set $\mathcal S\subset\mathcal F$ supplied by
Theorem~\ref{fin:obstruction}.  We shall use both its double-coset
obstruction and its assertion that, on an abelian subgroup meeting
$\mathcal S$, the complement of $\mathcal S$ is a subgroup.  In particular,
$\mathcal S$ is inverse-stable.  The probability construction and the
fractional transformations below adapt
\cite[Sections~5.2--5.4]{MPNF}.  Section~\ref{sec:recurrence} supplies the
additive argument in positive characteristic.

\subsection{A compact space of color configurations}

Let $\widetilde G$ be the special unitary group of the hyperbolic
Hermitian form on $D^2$ with coefficients $1,-1$, and let $X$ be its
projective variety of isotropic right $D$-lines.  This is a flag variety
for a minimal $k$-parabolic of $\widetilde G$: the division hypothesis
makes its relative $k$-rank one.  Every rational isotropic line is
represented uniquely by a column $(u,1)$ with $u\in U(k)$.  Indeed, its
second coordinate cannot be zero, and every nonzero element of $D$ is
invertible.  We henceforth identify $X(k)$ with $U(k)$ in this way.

Choose $\theta\in L\setminus k$, with $\bar\theta=-\theta$ when
$\operatorname{char}k\ne2$, and put
\[
 J=\{x\in D:x^*=x\},\qquad
 q(x)=(x+\theta)(x+\bar\theta)^{-1}.
\]
The space $J$ is an additive $k$-vector group; its rational points will
also be denoted by $J$.  The column $(x+\theta,x+\bar\theta)$ is isotropic,
and the Cayley parametrization gives
\begin{equation}\label{pal:cayley}
 X(k)=J\sqcup\{\infty\},\qquad q(\infty)=1.
\end{equation}
For example, $x+\bar\theta$ cannot vanish for Hermitian $x$, and the
inverse parametrization is defined at every $u\ne1$ because $u-1$ is
invertible.  Over a completion, the unitary and additive parametrizations
are open charts in $X(k_v)$; their complements are proper algebraic
subvarieties and have measure zero in its smooth measure class.

Let $\mathcal H$ be the countable group of rational transformations of
$X$ induced by form similitudes.  Besides $\widetilde G(k)$, it contains
\[
 L_a(u)=au,\quad R_a(u)=ua\quad(a\in U(k)),\qquad
 T_b(x)=x+b\quad(b\in J).
\]
It also contains the central affine maps $x\mapsto lx+s$, with
$l\in k^\times$ and $s\in k$.  We write $\tau_s=T_s$ for scalar
translations.  For $p\in X(k)$ define
\[
 c(p)=(c_h(p))_{h\in\mathcal H},\qquad
 c_h(p)=\phi(u(hp)).
\]
The closure of these configurations in $\mathcal F^{\mathcal H}$ is
denoted by $\mathcal C$; its members are called \emph{palettes}.  It is
a compact metrizable space.  The action and its relation to rational
points are
\[
 (g c)_h=c_{hg},\qquad c(gp)=g c(p).
\]
Every finite coordinate pattern occurring in $\mathcal C$ occurs at a
rational point.  Consequently identities involving finitely many colors
pass to all palettes.  For instance,
\begin{equation}\label{pal:rotation-rules}
 c_{L_a h}=\phi(a)c_h,\qquad
 c_{R_a h}=c_h\phi(a).
\end{equation}

\subsection{An invariant law retaining all colors}

A translation-invariant probability is easy to obtain by averaging.
What requires proof is that averaging can be arranged to retain every
color at every coordinate.  We first construct a joint probability on
local flags and palettes, whose conditional color distributions are
uniform.

\begin{proposition}\label{pal:uniform-law}
There is a probability measure $\mu$ on $\mathcal C$ such that
\[
 (T_b)_*\mu=\mu\quad(b\in J),\qquad
 \mu\{c:c_h=a\}=|\mathcal F|^{-1}
 \quad(h\in\mathcal H, a\in\mathcal F).
\]
\end{proposition}

We divide the proof into the construction of a joint law, identification
of its base measure, and uniformity of its conditional colors.

\subsubsection{Rational weights and a joint law}

For each place $v$, choose an Iwasawa compact subgroup
$\mathcal K_v\subset\widetilde G(k_v)$ and its invariant probability
$m_v$ on $X(k_v)$.  The compact group is transitive on this flag variety.
At almost all places use the integral models.  Set
\[
 Y=\prod_vX(k_v),\qquad m=\prod_vm_v.
\]
Each local probability belongs to the smooth measure class.  A rational
similitude preserves $m_v$ at all but finitely many places: away from its
denominators and the places where its multiplier is a nonunit, it
normalizes the chosen integral compact.  Thus
\[
 \rho(h,y)=\frac{d(h_*m)}{dm}(y)
\]
is a continuous positive function given by a finite product of local
densities.  With this pushforward convention its cocycle identity is
\begin{equation}\label{pal:cocycle}
 \rho(ag,y)=\rho(a,y)\rho(g,a^{-1}y).
\end{equation}

Assign a positive weight $w(p)$ to each rational point by transporting
weight $1$ from a fixed base point under $\widetilde G(k)$, multiplying
by the forward local volume Jacobians.  Rational conjugacy of parabolics
gives transitivity.  The weight is independent of the transporter:
at a rational fixed point the product of tangent determinants is $1$
by the product formula.  The same argument gives the weight rule for
every $h\in\mathcal H$, including similitudes.

We need convergence of these weights before they can define a probability.
The following proof uses the discreteness of the function-field modulus
and makes the exponent explicit.

\begin{lemma}\label{pal:weight-sum}
For every real $t>1$,
\[
 \sum_{p\in X(k)}w(p)^t<\infty.
\]
\end{lemma}

\begin{proof}
Let $P$ be the stabilizer of the base point, with unipotent radical $N_P$.
Write $\delta_P$ for its adelic modulus, the absolute value of the
determinant of its action on $\Lie(N_P)$.  Extend it to
$\widetilde G(\bbA_k)$ by Iwasawa decomposition with a right compact
factor:
\[
 g=p_0k_0,\qquad \delta_P(g)=\delta_P(p_0),\qquad
 k_0\in\mathcal K:=\prod_v\mathcal K_v.
\]
For a stabilizing element the forward tangent Jacobian is the inverse
modulus.  Taking the rational transporter to be $\gamma^{-1}$ therefore
identifies the required sum with
\begin{equation}\label{pal:eisenstein-sum}
 \sum_{\gamma\in P(k)\backslash\widetilde G(k)}\delta_P(\gamma)^t.
\end{equation}

The values $\delta_P(\gamma)$ in this sum are bounded above.  To see this,
take a rational vector $v_P$ spanning the top exterior power of
$\Lie(N_P)$, considered inside the corresponding exterior-power
representation of $\widetilde G$.  Choose local norms integral almost
everywhere.  The Iwasawa decomposition gives
\[
 \prod_v\|\gamma^{-1}v_P\|_v\leq C_1\delta_P(\gamma)^{-1}.
\]
The product norm of a nonzero rational vector is bounded below by a
fixed positive constant: one nonzero rational coordinate and the product
formula suffice.  Hence $\delta_P(\gamma)\leq C$ uniformly.

Choose a compact open subgroup $K'\subset\mathcal K$ so small that
$K'\cap\widetilde G(k)=\{1\}$.  The images of the sets $\gamma K'$ in
$P(k)\backslash\widetilde G(\bbA_k)$ are disjoint for distinct cosets in
\eqref{pal:eisenstein-sum}, and all have the same positive volume.
Indeed an overlap would put $\gamma_2^{-1}p\gamma_1\in\widetilde G(k)$
in $K'$ for some $p\in P(k)$.  The same argument excludes stabilizers
on each such neighborhood.  Since the modulus is right-$K'$-invariant,
it suffices to show that
\[
 \int_{P(k)\backslash\widetilde G(\bbA_k)}
 \mathbf1_{\{\delta_P\leq C\}}\delta_P(g)^t\,dg<\infty.
\]

Iwasawa integration reduces this to left Haar integration over
$P(k)\backslash P(\bbA_k)$.  One can normalize Haar measures by the
compact groups $\mathcal K$ and $P(\bbA_k)\cap\mathcal K$; the resulting
double cosets and their stabilizer weights agree.  Let $M_P$ be a Levi
subgroup, and use the decomposition $P=N_P M_P$, with the unipotent
coordinate first.  The left Haar measure contains the factor
$\delta_P(m_0)^{-1}$.  The quotient of the unipotent radical has fixed
finite volume, so integration along it leaves
\[
 \mathbf1_{\{\delta_P(m_0)\leq C\}}
 \delta_P(m_0)^{t-1}
\]
on $M_P(k)\backslash M_P(\bbA_k)$.  The product formula makes the
fiber volumes independent of rational changes of coordinates.

The group $\widetilde G$ has relative rank one and $P$ is minimal.
Thus the kernel of the modulus on the adelic Levi is its norm-one
adelic group, whose rational quotient has finite volume by reduction
theory; see \cite{HarderReduction}.  The nonempty modulus levels all
have this same finite volume, since Levi Haar measure is also right
invariant.  The modulus values form a discrete subgroup of the powers
of the cardinality of the full constant field.  Summing
$\delta_P^{t-1}$ over the levels bounded above is a convergent geometric
series precisely when $t>1$.  This proves the lemma.
\end{proof}

For $t>1$, give $(p,c(p))\in Y\times\mathcal C$ normalized mass
$w(p)^t$, and denote the resulting probability by $\mu_t$.
The weight rule gives
\[
 h_*\mu_t=\rho(h,\cdot)^t\mu_t.
\]
The reciprocal implicit in the pushforward density is accounted for
here: the mass at $hp$ after pushforward is the old mass at $p$.
By compactness take a weak limit $\mu^0$ as $t\downarrow1$.  Uniform
convergence of $\rho(h,\cdot)^t$ for each fixed $h$ gives
\begin{equation}\label{pal:joint-law}
 h_*\mu^0=\rho(h,\cdot)\mu^0\qquad(h\in\mathcal H).
\end{equation}

\subsubsection{The base projection is the product probability}

Let $\nu$ be the projection of $\mu^0$ to $Y$.  We show that $\nu=m$;
in particular, no singular base probability enters the later conditional
argument.  Fix a sufficiently large nonempty finite set $T$ of places.
Subscripts $T$ denote products over $T$, and superscripts $T$ denote
the complementary products.  An arithmetic lattice
$\Gamma\subset\widetilde G_T$ integral outside $T$ preserves $m^T$.
Consequently \eqref{pal:joint-law} gives
\[
 \gamma_*\nu_T=\rho_T(\gamma,\cdot)\nu_T\qquad(\gamma\in\Gamma).
\]
Define
\[
 f(g)=\int_{Y_T}\rho_T(g,y)\,d\nu_T(y).
\]
The cocycle rule and this transformation law imply
$f(\gamma g)=f(g)$.  The function is positive and continuous on
$\widetilde G_T$.

Choose a probability $\eta$ of full support on $\widetilde G_T$ that is
invariant under left multiplication by $\mathcal K_T$.  Averaging first
over that compact group shows that
\[
 \int\rho_T(h,z)\,d\eta(h)=1\qquad(z\in Y_T):
\]
compact transitivity turns the inner average into the integral of a
probability density over $m_T$.  By \eqref{pal:cocycle}, it follows that
\[
 f(g)=\int f(gh)\,d\eta(h).
\]
Apply Jensen's inequality to the bounded strictly concave function
$a\mapsto a/(1+a)$ on the positive ray.  Its nonnegative Jensen
discrepancy has integral zero on the finite-volume quotient
$\Gamma\backslash\widetilde G_T$, by right invariance of quotient
measure.  Equality in strict Jensen implies that $f(gh)$ is constant
for $\eta$-almost every $h$, for almost every $g$.  Full support and
continuity then make $f$ constant everywhere.  Its value at $1$ is $1$.
This argument requires no integrability assumption on $f$ itself.

These density integrals determine the base probability.  Choose
$a_j\in\widetilde G_T$ contracting a big cell to a point $x_0\in Y_T$,
using a parabolic cocharacter in each local factor.  The complement of
the cell has measure zero, so $(a_j)_*m_T\to\delta_{x_0}$.  For
$b\in C(Y_T)$ set
\[
 B_j(y)=\int_{\mathcal K_T}b(k_0x_0)\rho_T(k_0a_j,y)\,dk_0.
\]
For fixed $y$, the density in this integral makes $k_0^{-1}y$ have law
$(a_j)_*m_T$.  This follows by pushing compact Haar measure to $Y_T$
and using $\rho_T(k_0a_j,y)=\rho_T(a_j,k_0^{-1}y)$.
Compact equicontinuity therefore gives $B_j(y)\to b(y)$ uniformly.
Since $f(k_0a_j)=1$,
\[
 \int B_j\,d\nu_T
 =\int_{\mathcal K_T}b(k_0x_0)\,dk_0
 =\int b\,dm_T.
\]
It follows that $\nu_T=m_T$.  Exhausting the places proves $\nu=m$.

Disintegrate the joint law as
\begin{equation}\label{pal:disintegration}
 d\mu^0(y,c)=dQ_y(c)\,dm(y).
\end{equation}
The compact metrizable spaces admit regular conditional probabilities.
Equation~\eqref{pal:joint-law} and uniqueness of disintegration give
$Q_{hy}=h_*Q_y$ almost everywhere.  Since $\mathcal H$ is countable,
we may use all these identities simultaneously.

\subsubsection{Uniform conditional colors}

Let $\mathbf p(y)$ be the $c_1$-marginal of $Q_y$, regarded as a
probability vector on $\mathcal F$, and put
$N=\ker\phi\cap S(k)$.  The palette rules
\eqref{pal:rotation-rules} imply that $\mathbf p$ is invariant under
the separate left and right rotations by every element of $N$.
We first prove invariance under a whole local special unitary group.

Enlarge $T$ to contain a place $v_0$ at which $S$ is split, as well as
the required integral-model exceptions.  On the unitary chart over $T$,
the base measure is in the Haar class of $U_T$.  Disintegrate that class
along the smooth determinant map, and choose measurable representatives
$d_\lambda$ of its fibers.  Write
\[
 u_T=g d_\lambda,\qquad g\in S_T.
\]
The conditional measure class in $g$ is Haar measure.  For almost every
determinant $\lambda$, each component of
$\mathbf p(gd_\lambda,y^T)$ descends, by its left invariance, to
\begin{equation}\label{pal:lattice-space}
 Z_\lambda=\Gamma'\backslash(S_T\times Y^T),\qquad
 \Gamma'=N\cap\{a\in S(k):a\text{ is integral outside }T\}.
\end{equation}
The group $\Gamma'$ has finite index in an arithmetic lattice.
The quotient has finite measure, obtained from Haar measure on a
lattice fundamental domain times $m^T$; the integral left tail actions
preserve the latter.  The bounded components of $\mathbf p$ thus
belong to $L^2(Z_\lambda)$.

Right invariance by $a\in\Gamma'$ acts on this quotient through
\[
 (g,y^T)\longmapsto
 (g d_\lambda a_Td_\lambda^{-1},R_a y^T).
\]
The tail rotations commute with the left action and belong to a compact
group of measure-preserving transformations.  There are elements
$a_j\in\Gamma'$ escaping every compact set at $v_0$ while remaining
bounded at the other factors of $S_T$.  For example, use
Proposition~\ref{arith:closure} to find infinitely many elements of
$R\subseteq N$ in a fixed compact-open cylinder away from $v_0$.
Adelic discreteness forces an escaping subsequence at $v_0$.
After another subsequence the bounded components and compact tail
rotations converge.  Conjugating by the fixed $d_\lambda$ does not
change these properties.

Decompose $L^2(Z_\lambda)$ into the invariant subspace for the right
$S(k_{v_0})$ action and its orthogonal complement.  The local group is the split group $\SL_n(k_{v_0})$, with its finite
center, and local Kneser--Tits identifies it with its root-generated
group.  Hence Howe--Moore coefficient decay applies on
the complement in every characteristic; see
\cite{HoweMoore} and \cite[Theorem~4.19]{Ciobotaru}.  If $v$ is the
projection of a component of $\mathbf p$ to that complement, the
operators fixing it have the form $\pi(g_j)A_j$, where $g_j$ escapes
in $S(k_{v_0})$ and the commuting unitaries $A_j$ converge strongly to
a unitary $A$.  Therefore
\[
 \|v\|^2=\langle\pi(g_j)A_jv,v\rangle
 =\langle\pi(g_j)Av,v\rangle+o(1)\longrightarrow0.
\]
The local and compact actions used here are strongly continuous, as
one first checks on continuous functions with compact support and then
by density in $L^2$.  We conclude that $\mathbf p$ is invariant under
the full local right group $S(k_{v_0})$.

Lifting back and using Haar measure and Fubini, $\mathbf p$ consequently
depends on $u_{v_0}$ only through its determinant, up to null sets.
More explicitly, restrict at $v_0$ to the conull unitary chart and
consider the measurable projection
\[
 \pi:Y\longrightarrow\overline Y
   :=\det(U(k_{v_0}))\times\prod_{v\ne v_0}X(k_v),\qquad
 y\longmapsto(\det u_{v_0},(y_v)_{v\ne v_0}).
\]
Give $\overline Y$ the pushforward probability $\overline m$.
The local invariance just proved says that
$\mathbf p=\overline{\mathbf p}\circ\pi$ almost everywhere for a
bounded measurable function $\overline{\mathbf p}$ on this space.
Every rational left rotation from $S(k)$ acts trivially on the first
coordinate of $\overline Y$.

Take $g\in V_0$.  Proposition~\ref{arith:closure} supplies
$r_j\in R$ tending to $g$ in the restricted adelic product away from
$v_0$.  Their left actions on $\overline Y$ converge to that of $g$.
This convergence passes to bounded measurable functions: outside one
fixed finite set the actions preserve the factor probabilities, and
on the remaining local factors their densities are uniformly bounded.
The determinant factor of $\overline m$ is unchanged.  For a bounded continuous
cylinder function, convergence follows from local continuity and
bounded convergence; approximation in $L^1(\overline m)$ then proves
it for $\overline{\mathbf p}$.  Since every $r_j$ belongs to
$R\subseteq N$, invariance passes to the limit and gives
$\mathbf p(L_g y)=\mathbf p(y)$ for every $g\in V_0$, almost everywhere.

Conditional equivariance also gives
$\mathbf p(L_g y)=\phi(g)\mathbf p(y)$.  Since $\phi(V_0)=\mathcal F$,
the probability vector is uniform.  Equivariance identifies the
$c_h$-marginal at $y$ with the $c_1$-marginal at $hy$, so all conditional
coordinate marginals are uniform.

\begin{proof}[Completion of the proof of Proposition~\ref{pal:uniform-law}]
Let $\eta_0$ be the palette marginal of $\mu^0$.  The countable additive
group $J$ is an increasing union of finite additive subgroups $J_j$.
Average
\[
 \eta_j=|J_j|^{-1}\sum_{b\in J_j}(T_b)_*\eta_0.
\]
Every summand has uniform coordinate marginals, because a translation
merely changes the coordinate index.  Any weak limit $\mu$ has the same
marginals.  For each fixed $b$, all sufficiently late averages are
$T_b$-invariant, so the limit is invariant under every translation.
\end{proof}

\subsection{A translation obstructed by a change of membership}

We now connect this law to Theorem~\ref{fin:obstruction}.  Given
$w\in U(k)$ and a scalar affine map $x\mapsto lx+s$, its image in
unitary coordinates agrees, up to a central norm-one scalar, with
\begin{equation}\label{pal:affine-fraction}
 z=(w-\bar t)(1-tw)^{-1}.
\end{equation}
Here, putting
\[
 a=l\bar\theta-s-\theta,\qquad b=s-(l-1)\theta,
\]
one takes $\bar t=-b/a$ when $a\ne0$.  Direct substitution in
\eqref{pal:cayley} proves this formula.  Moreover
\[
 a\bar a-b\bar b=-l(\theta-\bar\theta)^2,\qquad
 d:=1-t\bar t\ne0.
\]
If $a=0$, then $(l,s)=(-1,-\theta-\bar\theta)$, and the affine
transformation inverts colors.  It therefore preserves membership in
$\mathcal S$ and will never cause an exceptional membership change.

For $a\ne0$ set
\[
 C=tw,\qquad R_C=(1-C^*)(1-C)^{-1},\qquad z=wR_C.
\]
These definitions agree with \eqref{pal:affine-fraction}.
The elements $C,C^*$ commute and $CC^*=t\bar t$ is central.
Division and $d\ne0$ make the denominators invertible; $R_C$ is
unitary and commutes with $w$.  In particular $z$ commutes with $w$.

Consider the similitude
\[
 D_C(u)=(u+C)(C^*u+1)^{-1}.
\]
Its matrix has multiplier $d$.  Its denominator cannot vanish at a
unitary element, since that would force $t\bar t=1$.
Writing $u=w\alpha$ gives the useful identity
\begin{equation}\label{pal:commuting-fraction}
 w^{-1}D_C(u)=(\alpha+t)(\bar t\alpha+1)^{-1},\qquad
 [w^{-1}D_C(u),w^{-1}u]=1.
\end{equation}
No commutation between $u$ and $w$ is being assumed.

For a variable $u\in U(k)$ put
\[
 v=D_{C^*}(uR_C),\qquad u'=R_C D_{-C}(v).
\]
Apply \eqref{pal:commuting-fraction} first with parameters
$w^{-1},\bar t$, and then with parameters $w,-t$.  Conjugating the
first resulting commutation gives
\begin{equation}\label{pal:two-commutations}
 [vw,uz]=1,\qquad [w^{-1}v,z^{-1}u']=1.
\end{equation}
On the other hand, matrix multiplication gives the fractional matrix
of $u\mapsto u'$ as
\[
 \begin{pmatrix}
 d+C-C^*&-(C-C^*)\\ C-C^*&d-(C-C^*)
 \end{pmatrix}.
\]
Applying it to $(x+\theta,x+\bar\theta)$ and dividing by $d$ shows
that it is the additive translation by
\begin{equation}\label{pal:translation-vector}
 b_w(l,s)=\frac{(\theta-\bar\theta)(C-C^*)}{d}\in J.
\end{equation}
The matrix calculation also applies at infinity.

To detect simultaneous returns, take one independent copy of
$(\mathcal C,\mu)$ for each color, and in the product space let $E_0$
be the event that the first coordinate of the copy indexed by
$a\in\mathcal F$ equals $a$.  Then
\[
 \delta:=\mathbb P(E_0)=|\mathcal F|^{-|\mathcal F|}>0.
\]
For the diagonal translation action define its return correlation
\begin{equation}\label{pal:correlation}
 K(b)=\mathbb P(E_0\cap T_b^{-1}E_0)\quad(b\in J).
\end{equation}
This is nonnegative and positive definite.  The spectral theorem for
the discrete abelian group $J$ gives a positive measure $\sigma$ on
its compact character group $\widehat J$ such that
\begin{equation}\label{pal:spectral-measure}
 K(b)=\int_{\widehat J}\chi(b)\,d\sigma(\chi),\qquad
 \sigma(\widehat J)=\delta,\qquad \sigma(\{1\})\geq\delta^2.
\end{equation}
The last inequality follows by projecting $\mathbf1_{E_0}$ onto
invariant vectors; its projection onto constants alone has squared
norm $\delta^2$.

\begin{proposition}\label{pal:translation-obstruction}
If the colors of $w$ and its scalar affine image satisfy
$\phi(w)\in\mathcal S$ and $\phi(z)\notin\mathcal S$, then
\[
 K(b_w(l,s))=0.
\]
\end{proposition}

\begin{proof}
If the correlation were positive, every color would occur as
$\phi(u)=\phi(u')$ for the fixed translation in
\eqref{pal:translation-vector}.  Indeed a point of the positive
intersection supplies the two-coordinate return pattern in each
palette copy, and every such finite pattern occurs rationally.
The rational points for different colors may differ, which is all
that is needed.

Write $W=\phi(w)$, $Z=\phi(z)$, and, for a return point, put
$X'=\phi(u)=\phi(u')$ and $Y'=\phi(v)$.  With $B=Z^{-1}X'$,
Equation~\eqref{pal:two-commutations} yields
\[
 W^{-1}Y'\in C_{\mathcal F}(B),\qquad
 Z^{-1}Y'WZ\in C_{\mathcal F}(BZ^2).
\]
As $WZ=ZW$, direct multiplication gives
\[
 (W^{-1}Y')^{-1}(ZW^{-1})(Z^{-1}Y'WZ)=ZW.
\]
The return colors make $B$ range over all of $\mathcal F$, contradicting
Theorem~\ref{fin:obstruction}.  The exceptional affine map preserves
membership and hence does not arise under the stated hypotheses.
\end{proof}

It remains to record the precise dependence of this translation on
$(l,s)$.  This will allow us to apply additive Fourier analysis to
scalar affine positions without treating multiplicative colors as
continuous functions.  Put
\[
 T=\theta+\bar\theta,\qquad P_\theta=\theta\bar\theta.
\]
For fixed $w$ define the $k$-linear map
\[
 \mathcal B_w(\alpha,\beta)=
 \frac{\alpha(w-w^*)+\beta(\bar\theta w-\theta w^*)}
      {\theta-\bar\theta}:k^2\longrightarrow J.
\]
Expanding $-a\bar b$ and using the preceding denominator identity gives
\[
 \frac{t}{d}=
 \frac{\alpha+\bar\theta\beta}{(\theta-\bar\theta)^2},\qquad
 \alpha=P_\theta l-\frac{s^2+Ts+P_\theta}{l},\quad
 \beta=2s+T(1-l).
\]
Consequently
\begin{equation}\label{pal:scalar-parameters}
 b_w(l,s)=\mathcal B_w\left(
 P_\theta l-\frac{s^2+Ts+P_\theta}{l},\,2s+T(1-l)\right).
\end{equation}
Both scalar arguments vanish at the exceptional affine parameter, so
this formula extends there by $b_w=0$.  Equations~\eqref{pal:spectral-measure}
and~\eqref{pal:scalar-parameters}, together with
Proposition~\ref{pal:translation-obstruction}, are the inputs to the
recurrence argument.
